\documentclass[11pt]{article}
\usepackage[margin=1in]{geometry}
\usepackage{amsmath,amssymb,amsthm,microtype}
\usepackage[hidelinks]{hyperref}
\newtheorem{theorem}{Theorem}
\newtheorem{lemma}[theorem]{Lemma}
\newtheorem{proposition}[theorem]{Proposition}
\newcommand{\dist}{\operatorname{dist}}
\newcommand{\TF}{\mathcal T}
\newcommand{\Cthree}{\mathcal C_3}
\newcommand{\Athree}{\mathcal A_3}
\newcommand{\PP}{\mathcal P}
\newcommand{\eps}{\varepsilon}
\title{A polynomial versus superpolynomial gap\\in vertex-distribution-free graph testing}
\author{Alec Kriebel\\{\small Independent researcher}\\
{\small\url{https://orcid.org/0009-0001-9320-500X}}}
\date{17 October 2026 \quad Version 1.0.0}
\begin{document}
\maketitle
\begin{abstract}
We give a negative answer to the comparison question of Gishboliner and
Shapira in Oberwolfach Report 19/2019, Problem~5, and their later
Problem~5.6. The fixed property consisting of $K_4$-free graphs that
become $3$-colorable after deleting at most one vertex has a polynomial
classical canonical tester, for every input size. In the
vertex-distribution-free model, every adaptive two-sided tester has
superpolynomial total sample-plus-adjacency cost as the proximity
parameter tends to zero. The lower bound follows by an elementary
tripartite-cover and heavy-apex reduction from the established two-sided
triangle-testing lower bound. All distances allow arbitrary edge
additions and deletions. The reduction remains valid when the tester is
given the vertex set, its size, the full positive rational distribution,
and the heavy vertex. A self-contained classical sampling argument and
exact finite verification materials accompany the note.
\end{abstract}

\section{The question and the fixed property}
All graphs are finite, simple and undirected. Sampling inputs have a
nonempty vertex set. For a probability distribution $D$ on that set, use
the unordered-pair normalization
\[
 d_D(G,F)=\sum_{\{u,v\}\in E(G)\triangle E(F)}D(u)D(v),
 \qquad
 \dist_D(G,\mathcal R)=\min_{F\in\mathcal R,\ V(F)=V(G)}d_D(G,F).
\]
Thus uniform distance is the number of changed edges divided by $n^2$.
We call distance at least $\eps$ $\eps$-far. Strict-threshold conventions
give the same conclusions after a fixed proximity rescaling.

A tester accepts members and rejects $\eps$-far inputs with probability
at least $2/3$. The VDF oracle supplies independent samples from the
unknown vertex distribution and adjacency answers. We count their
\emph{total} number, as in the original workshop model
\cite[p.~1138]{OWR}. The comparison question asks whether an extendable
hereditary property has the same, or nearly the same, testing complexity
in the VDF and uniform dense models \cite[Problem~5]{OWR},
\cite[Problem~5.6]{GS}. Hereditary means closed under induced subgraphs;
extendable means that every member has a one-vertex extension which is
also a member.

Let $\Cthree$ be $3$-colorability, let $\Athree$ mean that deletion of
at most one vertex leaves a $3$-colorable graph, and fix
\[
                    \PP=\{K_4\text{-free graphs}\}\cap\Athree .
\]
This property is hereditary, closed under deleting edges, and
extendable by adding an isolated vertex. Moreover, $\Cthree\subseteq\PP$.

\begin{theorem}\label{thm:main}
For this single fixed property $\PP$:
\begin{enumerate}
\item There is a one-sided classical canonical tester, without knowing
$n$, using $O(\eps^{-6}\log(1/\eps))$ independent uniform samples and
$O(\eps^{-12}\log^2(1/\eps))$ adjacency queries as $\eps\to0$.
\item Along a sequence of positive rational parameters $\eps\to0$,
every two-sided adaptive VDF tester has worst-case total oracle cost
$\exp(\Omega(\log^2(1/\eps)))$. This holds even with the vertex set,
its size, the complete positive rational weight vector and the
distinguished heavy vertex supplied.
\end{enumerate}
Consequently VDF complexity admits no fixed polynomial bound in the
classical canonical complexity and $1/\eps$, even after any fixed
positive constant rescaling of classical proximity.
\end{theorem}

The lower-bound ingredients are established results: polynomial
colorability sampling \cite{GGR,AK}, the hard-triangle constructions and
their two-sided extensions \cite{Alon,AS,AKKR}, and the VDF framework
\cite{Goldreich,GS}. The reduction below transfers triangle hardness to
a property with an easy classical canonical tester. It supplies the
full negative comparison, rather than a new triangle lower bound.
Fei's recent edge-distribution work \cite{Fei} also uses classical
parity constructions; its default sample oracle and arbitrary pair
weights differ from the present adjacency oracle and vertex-product
distance. The current priority audit, including primary theorem bodies
and version histories, identified no earlier full comparison answer.
This is a bounded literature finding: the 2021 journal version-of-record
body of \cite{GS} was unavailable, although its arXiv v5, ECCC revision~2
and author manuscripts were inspected; forward-citation retrieval was
incomplete. Detailed coverage and limitations accompany the package.

\section{A classical tester at every input size}
We include a conservative form of the standard branching-list
colorability argument \cite{GGR,AK}, so no large-input qualification is
needed.

\begin{lemma}\label{lem:colour}
For $0<\eps\le1/8$, put
\[
 h=\lfloor6/\eps\rfloor+1,\qquad
 \ell=\left\lceil(2/\eps)\log(12\cdot3^h)\right\rceil,\qquad r=h\ell .
\]
If a nonempty graph is $\eps$-far from $\Cthree$, then $r$ independent
uniform samples induce a graph outside $\Cthree$ with probability at
least $11/12$. Repeated labels are identified.
\end{lemma}
\begin{proof}
Fix a proper partial coloring on $S$. For $v\notin S$, let $L(v)$ be
the colors absent from its neighbors in $S$; let $U$ comprise the
empty-list vertices. For $v\notin S\cup U$, define
\[
 \delta(v)=\min_{i\in L(v)}
 |\{u\notin S\cup U: uv\in E(G),\ i\in L(u)\}|.
\]
Color every such $v$ by a minimizing color and color $U$ arbitrarily.
Deleting monochromatic edges gives a $3$-colorable graph with at most
$n|U|+\sum_v\delta(v)$ changes. If
$W=\{v\notin S\cup U:\delta(v)\ge\eps n/2\}$, farness implies
\[
 \eps n^2\le n|U|+\sum_v\delta(v)
       \le n(|U|+|W|)+\eps n^2/2 .
\]
Hence $|U\cup W|\ge\eps n/2$. The list potential
$\Phi=\sum_{v\notin S}|L(v)|$ starts at $3n$. Coloring a vertex of
$W$ with any feasible color reduces it by at least $\eps n/2$;
a vertex in $U$ has no feasible color.

Partition the samples into $h$ independent blocks of $\ell$ samples.
At each node of a ternary tree of proper partial colorings, choose the
first sample in that level's block lying in its $U\cup W$, and branch
on its feasible colors. A node fails if its block misses $U\cup W$.
Conditional on all previous blocks, this probability is at most
$\exp(-\eps\ell/2)\le(12\cdot3^h)^{-1}$. There are fewer than $3^h$
possible node addresses at depths $0,\ldots,h-1$. Summing these
conditional bounds shows that no node fails with probability at least
$11/12$; nodes sharing a block need not be independent.

On this event no feasible branch survives $h>6/\eps$ levels, since
$\Phi$ stays nonnegative. A proper coloring of the sampled induced
graph would follow a surviving branch, a contradiction. Each chosen
vertex lies outside its path's previously colored set, so repeated
samples cause no difficulty.
\end{proof}

Set $N=12r^2$ and $M=\lceil N\log(3N)\rceil$. Sample $M$ vertices
uniformly with replacement, query every pair of distinct sampled
labels, and accept exactly when the induced graph belongs to $\PP$.
Membership can be decided by finite exhaustive search; the assertion
concerns oracle complexity, with no polynomial running-time claim.
Completeness is one-sided by heredity.

If $G$ is $\eps$-far from $\PP$, it is also $\eps$-far from
$\Cthree$. If $n\ge N$, the first two $r$-blocks both induce
non-$3$-colorable graphs and have disjoint label sets with probability
at least
\[
              1-\frac1{12}-\frac1{12}-\frac{r^2}{n}\ge\frac34 .
\]
Deleting one vertex cannot destroy both disjoint witnesses, so the
whole sample lies outside $\Athree$, hence outside $\PP$.
If $1<n<N$, the probability of missing any vertex is at most
\[
 n(1-1/n)^M\le n e^{-M/n}\le N e^{-M/N}\le1/3 .
\]
The case $n=1$ is immediate. Seeing all vertices rejects every far
input. The algorithm need not know which case holds; $M\ge2r$.
Since $h=O(\eps^{-1})$, $\ell=O(\eps^{-2})$ and $r=O(\eps^{-3})$,
this proves part~1 of Theorem~\ref{thm:main}. For $\eps>1/8$, use
parameter $1/8$. The empty graph belongs to $\PP$ and can be accepted
if supplied as an explicit empty input; it carries no sampling
distribution.

\section{Transferring triangle hardness}
Write $\TF$ for triangle-freeness.

\begin{lemma}[Heavy apex]\label{lem:apex}
If $H$ is $3$-colorable on $n\ge1$ vertices, add a universal vertex $a$
to form $J(H)$ and set $D(a)=1/2$, $D(v)=1/(2n)$ on the tail. Then
\[
             \dist_D(J(H),\PP)=\tfrac14\dist_{\rm unif}(H,\TF).
\]
\end{lemma}
\begin{proof}
Delete a minimum triangle-hitting set of tail edges. The resulting
tail remains $3$-colorable and triangle-free, so its join with $a$ is
$K_4$-free and belongs to $\PP$. Each tail edge costs $1/(4n^2)$,
giving the upper bound.

For any repair $F\in\PP$, let $S$ be the tail vertices whose edge to
$a$ was deleted. Its apex-edge cost is $\delta_a=|S|/(4n)$; denote
its tail-edit cost by $\delta_t$, including additions. The restriction
$F[V(H)\setminus S]$ is triangle-free, since a triangle together with
$a$ would give a $K_4$. Use that restriction and isolate $S$ to repair
$H$. This changes at most $4n^2\delta_t+n|S|$ edges. Therefore
\[
 \dist_{\rm unif}(H,\TF)\le4\delta_t+|S|/n
                      =4(\delta_t+\delta_a).
\]
Minimizing over all repairs proves the reverse inequality.
\end{proof}

\begin{lemma}[Tripartite cover]\label{lem:cover}
For any graph $H$ on $n\ge1$ vertices, let $L(H)$ have vertex set
$\{1,2,3\}\times V(H)$, with $(i,u)(j,v)$ an edge exactly when
$i\ne j$ and $uv\in E(H)$. Then $L(H)$ is $3$-colorable,
$H\in\TF$ implies $L(H)\in\TF$, and
\[
              \dist_{\rm unif}(L(H),\TF)
                      \ge\tfrac19\dist_{\rm unif}(H,\TF).
\]
\end{lemma}
\begin{proof}
The three parts color $L(H)$. A triangle projects to three distinct
pairwise adjacent vertices of $H$. Triangle-free repairs can be
restricted to deletions: intersecting a repair's edge set with the
original graph preserves triangle-freeness and never increases cost.
Take a minimum deletion repair of $L(H)$, with $t$ deleted edges.
Delete $uv$ from $H$ if any of its six cross-copy edges was deleted.
This deletes at most $t$ original edges. Every surviving triangle
would lift, one vertex in each part, to a surviving triangle in the
cover, a contradiction. Dividing the two deletion counts by $n^2$
and $(3n)^2$ gives the claim.
\end{proof}

Combining the lemmas gives, for every $H$,
\[
 \dist_D(J(L(H)),\PP)\ge
                   \frac1{36}\dist_{\rm unif}(H,\TF),                 \tag{1}
\]
where $D(a)=1/2$ and every other vertex has weight $1/(6n)$.
Triangle-free $H$ gives a member of $\PP$.

To use a VDF tester at proximity $\eps$, simulate it on
$(J(L(H)),D)$. A sample is generated locally from this explicit
distribution. Self-queries return zero, distinct apex--tail pairs are adjacent, and
same-part pairs or copies of the same original vertex are nonadjacent.
Every remaining adjacency uses at most one query to $H$.
All promised extra information is supplied explicitly. Consequently
a VDF tester with total budget $Q$ gives a two-sided adaptive tester
for triangle-freeness at proximity $36\eps$, using at most $Q$
adjacency queries. Samples require no graph queries. If the VDF execution uses $s$ samples
and $e$ adjacency calls, the sampled original indices and query endpoints
involve at most $s+2e\le2Q$ source vertices. This simulation
also works for randomized, history-dependent choices of queries and
for early stopping.

The established two-sided triangle-testing lower bound is
$\exp(\Omega(\log^2(1/\delta)))$ in the uniform dense graph model:
see \cite[Theorem~4(ii) and \S7]{AS}, including its explicit
undirected consequence. That work gives the vertex count and permits
adaptive adjacency queries; its inspected-vertex convention, with all
induced pairs available, changes the present query count by at most a
constant factor. Fixed changes in the unordered/ordered distance
normalization also leave this bound unchanged. Applying it at
$\delta=36\eps$ to the simulation proves part~2. It suffices to retain
a rational sequence of parameters; the weights above are already
strictly positive and rational. Combining this lower bound with the
polynomial classical bound proves the final assertion of
Theorem~\ref{thm:main}.

\section{Scope and verification}
The separation concerns total VDF oracle cost. A sample-only lower
bound with unlimited free arbitrary-name adjacency queries is not
asserted. The original proof preserved in the supporting package also
gives an explicit rational hard family and a direct parity/exposure
proof, with all its combinatorial and probabilistic estimates.
The shorter proof here imports the established triangle lower bound
and gives the complete reduction instead.

Independent programs test arbitrary weighted repairs on small graphs,
the cover reduction, coloring lists and finite adaptive queries.
These finite checks test the arguments; the universal claims rest on
the proofs and the cited triangle theorem. The original programs and their hashes accompany the reports.
The public checkers raise errors explicitly, and run their checks
with Python optimization enabled as well as in ordinary mode.

\paragraph{AI use and review status.}
AI tools were used extensively in developing and checking the
arguments, computations, literature review and writing. The note and
supplements are unrefereed. Independent adversarial reviews were
AI-assisted; no external human peer review, journal acceptance or
formal proof-assistant certification is claimed.

\begin{thebibliography}{10}\small
\bibitem{OWR}
L.~Gishboliner, joint work with A.~Shapira,
contribution to \emph{Combinatorics, Probability and Computing},
Oberwolfach Report 19/2019, Problem~5, p.~1139. 
\href{https://ems.press/content/serial-article-files/46798}{Publisher report}.

\bibitem{GS}
L.~Gishboliner and A.~Shapira,
\emph{Testing graphs against an unknown distribution},
Israel J. Math. \textbf{245} (2021), 787--837,
\href{https://doi.org/10.1007/s11856-021-2228-8}{doi:10.1007/s11856-021-2228-8}.
Inspected manuscript:
\href{https://arxiv.org/abs/1905.09903v5}{arXiv:1905.09903v5}, \S5.3.

\bibitem{Goldreich}
O.~Goldreich, \emph{Testing graphs in vertex-distribution-free models},
STOC 2019, 527--534. Inspected extended author manuscript,
4 October 2021.
\url{https://www.wisdom.weizmann.ac.il/~oded/R3/vdf.pdf}.

\bibitem{GGR}
O.~Goldreich, S.~Goldwasser and D.~Ron,
\emph{Property testing and its connection to learning and approximation},
J. ACM \textbf{45} (1998), 653--750.

\bibitem{AK}
N.~Alon and M.~Krivelevich, \emph{Testing $k$-colorability},
SIAM J. Discrete Math. \textbf{15} (2002), 211--227.

\bibitem{Alon}
N.~Alon, \emph{Testing subgraphs in large graphs},
Random Structures Algorithms \textbf{21} (2002), 359--370.

\bibitem{AS}
N.~Alon and A.~Shapira, \emph{Testing subgraphs in directed graphs},
J. Comput. System Sci. \textbf{69} (2004), 354--382,
\href{https://doi.org/10.1016/j.jcss.2004.04.008}{doi:10.1016/j.jcss.2004.04.008}.
Inspected author manuscript:
\url{https://www.math.tau.ac.il/~asafico/subdir.pdf}.

\bibitem{AKKR}
N.~Alon, T.~Kaufman, M.~Krivelevich and D.~Ron,
\emph{Testing triangle-freeness in general graphs},
SIAM J. Discrete Math. \textbf{22} (2008), 786--819.

\bibitem{Fei}
Y.~Fei, \emph{Testing properties of edge distributions},
\href{https://arxiv.org/abs/2603.22702v2}{arXiv:2603.22702v2},
12 August 2026, especially \S7.1.
\end{thebibliography}
\end{document}
